\documentclass[english]{uw-thesis}

\titleEN{Title of the doctoral dissertation}
\titlePL{Polski tytuł rozprawy doktorskiej}

\author{Your Full Name}
%\place{Warsaw}
\date{June 2026}

\faculty{Doctoral School of Exact and Natural Sciences}
%\faculty{International Doctoral School}
%\faculty{Faculty of Mathematics, Informatics, and Mechanics}

\addSupervisor{Prof.~Main Supervisor}
%\addSupervisor{Prof.~Second Supervisor}
\addCosupervisor{Dr~Assistant Supervisor}
\addCosupervisor{Dr~Second Assistant}
\addCosupervisor{Dr~Third Assistant}

%\fieldEN{Natural Sciences}
%\fieldPL{nauk ścisłych i przyrodniczych}

\disciplineEN{Computer and Information Sciences}
%\disciplineEN{Mathematics}
\disciplinePL{informatyka}
%\disciplinePL{matematyka}

\keywords{keyword one; keyword two; keyword three}

\abstractPL{Tutaj wstaw streszczenie po polsku.}
\abstractEN{Put the English abstract here.}

\newtheorem{theorem}{Theorem}[chapter]
\newtheorem{lemma}[theorem]{Lemma}

\begin{document}
\makefrontmatter

\chapter{Introduction}
Start the thesis here.

\section{First section}

\begin{theorem}
This is a sample theorem. You can use the \texttt{theorem} environment for theorems, lemmas, propositions, etc.
\end{theorem}

\cite{tarski_decision}

\section{Second section}

\begin{theorem}
This is a second theorem!
\end{theorem}

\begin{lemma}
This is a sample lemma.
\end{lemma}

\chapter{Contents}

This is the second chapter.

\begin{lemma}
This is a sample lemma.
\end{lemma}

\begin{lemma}
This is a sample lemma.
\end{lemma}

\bibliographystyle{alpha}
\bibliography{test.bib}
\end{document}
